\uuid{tR3p}
\titre{Calcul d'intégrales — puissances et fractions rationnelles}
\niveau{L1}
\module{Analyse}
\chapitre{Calcul d'intégrales}
\sousChapitre{Fraction rationnelle}
\theme{Décomposition en éléments simples, primitives de puissances}
\auteur{}
\datecreate{2026-06-04}
\organisation{}
\difficulte{2}

\contenu{
\texte{
  Calculer les intégrales suivantes.
}
\begin{enumerate}
\item
  \question{Calculer $\displaystyle\int_0^1 (1-x)\sqrt{x}\;dx$.}
  \reponse{
    \[
      \int_0^1 (1-x)\sqrt{x}\;dx
      = \int_0^1 \left(x^{1/2} - x^{3/2}\right)dx
      = \left[\frac{2}{3}\,x^{3/2} - \frac{2}{5}\,x^{5/2}\right]_0^1
      = \frac{2}{3} - \frac{2}{5}
      = \frac{4}{15}.
    \]
  }

\item
  \question{Calculer $\displaystyle\int_2^3 \dfrac{2x^2-1}{x(x-1)}\;dx$.
    On pourra déterminer $a$, $b$ et $c$ tels que
    $\dfrac{2x^2-1}{x(x-1)} = a + \dfrac{b}{x} + \dfrac{c}{x-1}$.}
  \indication{Effectuer la division euclidienne de $2x^2 - 1$ par $x^2 - x$,
    puis décomposer le reste en éléments simples.}
  \reponse{
    On effectue la division : $2x^2 - 1 = 2(x^2 - x) + (2x - 1)$, donc
    \[
      \frac{2x^2-1}{x(x-1)} = 2 + \frac{2x-1}{x(x-1)}.
    \]
    On pose $\dfrac{2x-1}{x(x-1)} = \dfrac{b}{x} + \dfrac{c}{x-1}$.
    En identifiant : $x = 0$ donne $b = 1$ ; $x = 1$ donne $c = 1$.
    Ainsi $a = 2$, $b = 1$, $c = 1$, et
    \begin{align*}
      \int_2^3 \frac{2x^2-1}{x(x-1)}\;dx
      &= \int_2^3 \left(2 + \frac{1}{x} + \frac{1}{x-1}\right)dx \\
      &= \Bigl[2x + \ln x + \ln(x-1)\Bigr]_2^3 \\
      &= (6 + \ln 3 + \ln 2) - (4 + \ln 2 + \ln 1) \\
      &= 2 + \ln 3.
    \end{align*}
  }
\end{enumerate}
}
